\section*{Bab \ref{ch:b3:cancer-biology} --- Biologi Kanker}

\begin{solution}{exo:b3:cancer-biology:1}
Sebuah \omterm{def:b3:cancer-biology:genes}{onkogen} adalah versi berfungsi-lebih sebuah gen yang memacu
proliferasi atau ketahanan hidup; satu alel mutan sudah memadai (dominan):
\emph{RAS} (mutasi titik), \emph{MYC} (translokasi, penguatan),
\emph{HER2}, \emph{BCR--ABL}. Sebuah penekan \omterm{def:b3:cancer-biology:hallmarks}{tumor} mengekang
proliferasi atau memaksakan penghentian, kematian, maupun perbaikan; kedua
alelnya harus hilang (resesif pada tingkat sel): \emph{RB}, \emph{TP53},
\emph{APC}, \emph{BRCA1}.
\end{solution}

\begin{solution}{exo:b3:cancer-biology:2}
Proliferasi yang berkelanjutan: \emph{RAS}, \emph{MYC}. Ketakpekaan terhadap
penekan pertumbuhan: \emph{RB}, \emph{CDKN2A} (p16). Menahan kematian:
\emph{BCL2}, \emph{TP53}. Kebakaan replikasi: promotor \omterm{def:b3:cancer-biology:telomerase}{telomerase}.
Angiogenesis: \emph{VEGF} yang diimbas lewat HIF (hilangnya \emph{VHL}).
\omterm{def:b3:cancer-biology:metastasis}{Invasi} dan \omterm{def:b3:cancer-biology:hallmarks}{metastasis}: hilangnya E-kaderin, serta faktor
epitelium--mesenkimnya. Ketakmantapan genom: gen \omterm{prop:b3:dna-repair:fidelity}{perbaikan salah
pasang}, \emph{\omterm{def:b3:dna-repair:dsb}{BRCA}}. Metabolisme yang tak teratur: \emph{MYC}, \omterm{prop:b3:cancer-biology:pathways}{jalur PI3K}.
\omterm{prop:b3:cancer-biology:immune}{Pengelakan kekebalan}: \emph{PD-L1}, hilangnya MHC. Peradangan yang memacu
\omterm{def:b3:cancer-biology:hallmarks}{tumor}: pengisyaratan \emph{NF-$\kappa$B}.
\end{solution}

\begin{solution}{exo:b3:cancer-biology:3}
Anak yang memiliki riwayat keluarga mengembangkan beberapa \omterm{def:b3:cancer-biology:hallmarks}{tumor} pada kedua
matanya, sejak dini; kasus sporadisnya mengembangkan satu \omterm{def:b3:cancer-biology:hallmarks}{tumor}, lebih lambat.
Ia menyimpulkan bahwa dua peristiwa dibutuhkan di dalam sebuah sel: pasien
keluarganya telah mewarisi satu lalu hanya menuntut peristiwa somatik yang
kedua (yang cukup sering terjadi pada beberapa sel retinanya), sedangkan
pasien sporadisnya menuntut keduanya terjadi secara somatik di dalam satu sel
--- langka, lambat, tunggal. Jadi sebuah gen yang kedua salinannya harus
hilang: yakni penekan \omterm{def:b3:cancer-biology:hallmarks}{tumor} yang pertama.
\end{solution}

\begin{solution}{exo:b3:cancer-biology:4}
Sebuah mutan \emph{Salmonella} yang menuntut histidina ditebarkan tanpa
histidina bersama senyawa ujinya; koloni hanya timbul dari sel yang
sebuah mutasi baru telah membalikkan cacatnya, sehingga cacahnya mengukur
daya mutageniknya. Ekstrak hati ditambahkan karena banyak karsinogen tak
berbahaya sampai enzim hati mengubahnya menjadi metabolit yang reaktif,
sebagaimana terjadi di dalam tubuh; bakterinya tak memiliki enzim itu.
\end{solution}

\begin{solution}{exo:b3:cancer-biology:5}
Insidensinya naik $100$ kali lipat sementara umurnya berlipat dua: $2^{k-1} = 100$, $k - 1 =
\log_{2}100 = 6.6$, sehingga kira-kira tujuh atau delapan peristiwa pembatas
laju pada pembacaan yang naif --- lebih sedikit bila pengembangan klonal di
antara pukulannya diperbolehkan.
\end{solution}

\begin{solution}{exo:b3:cancer-biology:6}
$L = \sqrt{2Dc_{0}/q}$: bagi $q = 0.01$, $\sqrt{2\times 2\times 10^{-9}
\times 0.05/0.01} = \qty{141}{\micro m}$; bagi $q = 0.03$,
\qty{82}{\micro m}. Sebuah \omterm{def:b3:cancer-biology:hallmarks}{tumor} yang tumbuh cepat bernapas lebih cepat,
menguras oksigennya dalam jarak yang lebih pendek, dan pusatnya mati pada
ukuran yang lebih kecil.
\end{solution}

\begin{solution}{exo:b3:cancer-biology:7}
$10^{11}\times 10^{-3n} < 1$ menuntut $n > 3.7$: empat penggal. Sesudah dua
penggal tersisa $10^{5}$ sel --- yakni sepersepuluh milimeter kubik, jauh
di bawah $10^{9}$ yang dapat dideteksi --- sehingga \omterm{def:b3:cancer-biology:hallmarks}{tumornya}
``lenyap'' menurut tiap ujinya padahal seratus ribu sel masih ada; pengobatan
harus dilanjutkan sampai cacah penggal yang direncanakan, bukan sampai \omterm{def:b3:cancer-biology:hallmarks}{tumornya}
menghilang.
\end{solution}

\begin{solution}{exo:b3:cancer-biology:8}
Kedua cacat itu mengaktifkan jalur yang sama (reseptor $\to$ \omterm{def:b3:cancer-biology:genes}{Ras} $\to$ ERK);
begitu salah satunya ada, yang lain tak memberi keunggulan lagi dan tak
terpilih. Sebuah penghambat EGFR menghalangi reseptornya; sebuah sel yang
memperoleh mutasi KRAS di hilirnya memulihkan isyaratnya dengan reseptornya
masih terhalang, lalu terpilih selama pengobatannya --- yakni kekebalan lewat
jalan pintas.
\end{solution}

\begin{solution}{exo:b3:cancer-biology:9}
E7 membebaskan E2F dari Rb sehingga mendorong sel epitelium yang terinfeksi
melewati \omterm{def:b3:cell-cycle-apoptosis:checkpoints}{titik batasnya} tanpa faktor pertumbuhan; E6 memusnahkan \omterm{def:b3:cancer-biology:genes}{p53}, sehingga
proliferasi yang tak pantas itu dan kerusakan yang ditimbulkannya tidak memicu
penghentian maupun \omterm{def:b3:cell-cycle-apoptosis:apoptosis}{apoptosis}. Dua ciri khasnya dipasok sekaligus, tetapi ciri
selebihnya --- kebakaan, \omterm{def:b3:cancer-biology:metastasis}{invasi}, ketakmantapan --- menuntut mutasi somatik
lanjutan, yang memakan beberapa dasawarsa untuk menumpuk di dalam sel yang
terinfeksi terus-menerus. Vaksinnya membangkitkan antibodi penetral terhadap
kapsidnya lalu mencegah infeksinya; bila diberikan sebelum paparan, tak pernah
ada sel terinfeksi yang dapat diubah.
\end{solution}

\begin{solution}{exo:b3:cancer-biology:10}
Sesudah pukulan ke-$i$, klonnya tumbuh sebesar $g$, sehingga $g^{i}$ kali
lebih banyak sel berisiko terkena pukulan berikutnya; peluang seluruh
runtunannya pada garis keturunan sebuah sel asal adalah $g\cdot g^{2}\cdots$
--- bagi $g$ yang tetap tiap pukulan, risiko kumulatifnya menjadi $N(ut)^{k}
g^{k-1}$ sampai faktor yang bergantung urutan, dan insidensinya mempertahankan
pangkat $t^{k-1}$ dengan prafaktor $g^{k-1}$ yang lebih besar. Bila
pengembangannya sendiri tumbuh dengan waktu (klon yang tumbuh eksponensial),
kurvanya mencuram lebih lanjut, sehingga kemiringan $6$ dihasilkan pemacu yang
kurang dari tujuh: $k$ Armitage--Doll itu batas atas cacah pemacu
pembatas laju, dan \omterm{def:b3:cancer-biology:hallmarks}{tumor} yang telah diurutkan memperlihatkan dua sampai delapan.
\end{solution}

\begin{solution}{exo:b3:cancer-biology:11}
(a) Obat A sendirian menghadapi $\mu_{A}N = 100$ sel kebal yang sudah ada:
$P_{0} = e^{-100} \approx 0$; klon yang kebal A bertahan lalu
tumbuh kembali selama delapan belas minggu pemberian A, dan pada saat B mulai
klon itu kembali menjadi populasi yang besar, sehingga B pun menghadapi
$\mu_{B}N' \gg 1$ --- runtunannya gagal dua kali. (b) Bersama: sebuah sel
kebal ganda timbul pada $10^{-14}$ tiap pembelahan, $\mu_{A}\mu_{B}N = 10^{-5}$, $P_{0} =
0.99999$. Gabungan sejak awal, ketika $N$ paling kecil, menjadi
kelaziman karena hasil kali dua laju yang kecil itulah yang membuat kekebalan
menjadi tidak mungkin.
\end{solution}

\begin{solution}{exo:b3:cancer-biology:12}
Dengan sel $100$ kali lebih banyak dan $u$ serta $k$ yang sama, risiko
kumulatif $N(ut)^{k}$ akan $100$ kali lebih tinggi --- setiap gajah semestinya
mati muda karena \omterm{def:b3:cancer-biology:hallmarks}{kanker}. Kenyataannya tidak, sehingga $u$ atau $k$ pasti
berbeda: dua puluh salinan \emph{TP53} membuat tanggapan apoptotik atas
kerusakan lebih kuat lalu menyingkirkan sel mutannya sebelum mengembang
(sehingga menurunkan $u$ efektif bagi langkah itu); hewan besar juga berlaju
mutasi tiap sel yang lebih rendah (perbaikan yang lebih baik, pembelahan tiap
sel tiap tahun yang lebih sedikit karena selnya berganti lebih lambat), dan
mungkin menuntut lebih banyak pukulan (sebuah lapis penekan tambahan).
Penekanan \omterm{def:b3:cancer-biology:hallmarks}{kanker} berevolusi bersama ukuran tubuh.
\end{solution}

\begin{solution}{pb:b3:cancer-biology:1}
\textbf{1.} $10^{12}/10^{3} = 10^{9}$ sel; $\log_{2}10^{9} \approx 30$
pelipatduaan.
\textbf{2.} $30\times 100 = 3000$ hari, kira-kira $8$ tahun.
\textbf{3.} \qty{1}{kg} $\approx 10^{12}$ sel: sepuluh pelipatduaan lagi,
$1000$ hari, kurang dari tiga tahun --- yakni sepertiga waktu sampai terdeteksi.
\textbf{4.} Waktu pelipatduaan awalnya pasti kira-kira $60$ hari, atau
\omterm{def:b3:cancer-biology:hallmarks}{tumornya} tumbuh lebih cepat ketika kecil lalu melambat seiring membesar (yakni
pola yang lazim, ketika pasokan dan ruangnya gagal).
\textbf{5.} Daur $2$ hari akan melipatduakan populasinya tiap $2$
hari; waktu pelipatduaan $100$ berarti pertumbuhan netonya $2\,\%$ laju
kelahirannya: nyaris tiap sel yang lahir diimbangi satu sel yang mati, atau
hanya sedikit yang berdaur --- kelahiran dan kematiannya nyaris berimbang, dan
pertumbuhannya selisih yang kecil.
\textbf{6.} Obatnya membunuh sel yang berdaur, sehingga \omterm{def:b3:cancer-biology:hallmarks}{tumor} dengan bagian
kecil yang berdaur kehilangan bagian yang lebih kecil tiap penggal; tetapi
\omterm{def:b3:cancer-biology:hallmarks}{tumor} yang cepat, yang telah kehilangan lebih banyak, juga tumbuh kembali di
antara penggalnya. Pada keseimbangannya, \omterm{def:b3:cancer-biology:hallmarks}{tumor} yang cepat (leukemia, limfoma,
\omterm{def:b3:cancer-biology:hallmarks}{kanker} buah zakar) itulah yang disembuhkan kemoterapi.
\textbf{7.} Nisbahnya $300$ atas nisbah umur $80/30 = 2.67$: $k - 1 =
\ln 300/\ln 2.67 = 5.8$, $k \approx 7$.
\textbf{8.} $N(ut)^{k} = 10^{8}\times (7\times 10^{-5})^{7} \approx
10^{-21}$: mustahil terhadap risiko seumur hidup yang mendekati $1$ dari $3$.
Yang kurang: pengembangan klonal yang melipatgandakan sel yang berisiko
sesudah tiap pukulan, laju mutasi yang dinaikkan ketakmantapan, pembelahan
yang jauh lebih banyak (sel punca yang berdaur berdasawarsa), serta pemacu
sungguhan yang lebih sedikit.
\textbf{9.} Mengalikannya dengan $g^{k-1} = 100^{6} = 10^{12}$ memberi $10^{-9}$
--- masih sangat kecil dengan tujuh peristiwa; dengan empat peristiwa dan
pengembangannya, $10^{8}\times (7\times 10^{-5})^{4}\times 10^{6} \approx 2\times
10^{-3}$, yakni orde yang tepat. Gambaran yang sebenarnya adalah sedikit
pemacu dan pengembangan yang besar.
\textbf{10.} $(ut)^{k-1}/(ut)^{k} = 1/(ut) = 1/(10^{-6}\times 40) =
2.5\times 10^{4}$: insidensi seorang pembawa pada umur empat puluh adalah
puluhan ribu kali insidensi sporadisnya --- \omterm{def:b3:cancer-biology:hallmarks}{kanker} turunan menyerang muda.
\textbf{11.} Satu peristiwa sepuluh kali lebih cepat: insidensinya $\times 10$;
dua peristiwa: $\times 100$. Risiko teramati sebesar lima belas sampai dua
puluh lima kali lipat menyarankan asapnya bekerja pada lebih dari satu langkah.
\textbf{12.} Begitu berhenti, $u$ turun kembali bagi peristiwa yang belum
datang, sehingga insidensinya berhenti mendaki setajam itu; tetapi pukulan
yang telah tertimbun di dalam selnya tetap ada, sehingga sel bekas perokok
lebih dekat ke \omterm{def:b3:cancer-biology:hallmarks}{kanker} daripada sel orang yang tak pernah merokok, dan risiko
lebihnya bertahan, membeku alih-alih terhapus.
\textbf{13.} $L = \sqrt{2\times 2\times 10^{-9}\times 0.05/0.03} =
\qty{82}{\micro m}$.
\textbf{14.} Beroksigen sepenuhnya sampai garis tengah $2L \approx
\qty{160}{\micro m}$. Sebuah sferoid \qty{2}{mm} bertepi hidup setebal
\qty{82}{\micro m}: bagian yang hidup $1 - (918/1000)^{3} = 0.23$.
\textbf{15.} Sel di tengahnya berjarak \qty{75}{\micro m} dari sebuah
pembuluh, tepat di dalam $L$ --- beroksigen secara asas, tetapi pembuluh \omterm{def:b3:cancer-biology:hallmarks}{tumor}
mengusung oksigen lebih sedikit dan mengalir terputus-putus, sehingga sel di
tengahnya hipoksia. Sel hipoksia dua sampai tiga kali lebih kebal terhadap
radiasi (oksigen dibutuhkan untuk membuat kerusakannya permanen), dan
hipoksia memicu \omterm{def:b3:cell-cycle-apoptosis:apoptosis}{apoptosis} yang bergantung \omterm{def:b3:cancer-biology:genes}{p53}, sehingga hipoksia memilih sel
yang telah kehilangan \omterm{def:b3:cancer-biology:genes}{p53}.
\textbf{16.} Memparuhkan kerapatannya merenggangkan pembuluhnya sebesar
$\sqrt{2}$: berjarak \qty{212}{\micro m}, sehingga sel di tengahnya
\qty{106}{\micro m} jauhnya, di seberang $L$: sel itu mati atau, bila bertahan
hidup dalam keadaan hipoksia, menjadi lebih \omterm{def:b3:cancer-biology:hallmarks}{ganas} dan lebih angiogenik.
\textbf{17.} $30/2 = 15$ kali lebih banyak glukosa. \omterm{def:b3:cancer-biology:hallmarks}{Tumor} karena itu
memekatkan analog glukosa, dan pemindaian fluorodeoksiglukosa yang memancarkan
positron menyalakannya.
\textbf{18.} \omterm{def:b3:cancer-biology:hallmarks}{Tumor} yang dilaparkan menciut sampai ukuran \omterm{prop:b3:cancer-biology:oxygen}{batas difusinya} lalu
bertahan, tertidur dan hipoksia, atau membajak pembuluh yang sudah ada;
hipoksianya memilih klon yang \omterm{def:b3:cancer-biology:hallmarks}{ganas}. Namun memangkas pembuluh \omterm{def:b3:cancer-biology:hallmarks}{tumornya} yang
kacau menurunkan tekanan antarselnya dan meratakan alirannya, sehingga
kemoterapinya menembus lebih baik --- gabungannya berhasil di tempat masing-masing
sendirian tidak.
\textbf{19.} $10^{9}\times 10^{-2n} < 1$: $n > 4.5$, lima penggal.
\textbf{20.} $21$ hari $= 0.21$ pelipatduaan, pertumbuhan kembalinya
$\times 1.16$; faktor neto tiap daurnya $0.0116$; $\ln 10^{-9}/\ln 0.0116 = 4.65$:
tetap lima daur.
\textbf{21.} $10^{-7}\times 10^{9} = 100$ sel kebal diharapkan;
$P_{0} = e^{-100} \approx 0$: kekebalannya pasti ada.
\textbf{22.} $\mu_{1}\mu_{2}N = 10^{-14}\times 10^{9} = 10^{-5}$;
$P_{0} = 0.99999$. Pada $10^{12}$ sel, $10^{-2}$, $P_{0} = 0.99$.
\textbf{23.} $10^{6}$ sel: satu obat $\mu N = 0.1$, $P_{0} = 0.90$; dua
obat $10^{-8}$, $P_{0} \approx 1$. Obatilah ketika populasinya kecil,
dan obatilah dengan gabungan.
\textbf{24.} Sel~T-nya menyerang banyak neoantigen sekaligus, sehingga tak
ada satu mutasi pun yang lolos darinya; pelolosannya menuntut hilangnya
penyajian antigennya sendiri (MHC, $\beta_{2}$-mikroglobulin) atau penekanan
sel~T-nya, yakni golongan peristiwa yang berbeda dan lebih langka. Muatan
mutasi yang tinggi yang sama, yang membuat sebuah \omterm{def:b3:cancer-biology:hallmarks}{tumor} pasti kebal terhadap
penghambat kinase, justru memberinya banyak neoantigen: \omterm{def:b3:cancer-biology:hallmarks}{tumor} yang \omterm{prop:b3:dna-repair:fidelity}{perbaikan
salah pasangnya} cacat serta \omterm{def:b3:cancer-biology:hallmarks}{tumor} imbasan ultraungu dan asap paling baik
menanggapinya.
\textbf{25.} Kira-kira $8$ tahun sampai terdeteksi; \omterm{def:b3:cancer-biology:hallmarks}{tumor} yang hidup setebal
\qty{82}{\micro m} di sekeliling sebuah pembuluh; $P_{0} = 0.99999$ bahwa
gabungan dua obat tak menghadapi satu pun sel kebal di dalam \omterm{def:b3:cancer-biology:hallmarks}{tumor} \qty{1}{cm^3}.
\end{solution}
